\documentclass[11pt,reqno]{amsart}
\usepackage[localtheorems]{raeez-math-template}

\usepackage{array}
\newtheorem{theorem}{Theorem}[section]
\newtheorem{proposition}[theorem]{Proposition}
\newtheorem{lemma}[theorem]{Lemma}
\newtheorem{corollary}[theorem]{Corollary}
\theoremstyle{definition}
\newtheorem{definition}[theorem]{Definition}
\newtheorem{example}[theorem]{Example}
\theoremstyle{remark}
\newtheorem{remark}[theorem]{Remark}
\newcommand{\cA}{\mathcal{A}}
\newcommand{\cB}{\mathcal{B}}
\newcommand{\cC}{\mathcal{C}}
\newcommand{\barB}{\overline{B}}
\newcommand{\barC}{\overline{C}}
\newcommand{\Conf}{\mathrm{Conf}}
\newcommand{\Conv}{\mathrm{Conv}}
\newcommand{\MC}{\mathrm{MC}}
\newcommand{\Ainf}{A_\infty}
\newcommand{\Vir}{\mathrm{Vir}}
\newcommand{\fg}{\mathfrak{g}}
\newcommand{\vac}{\mathrm{vac}}
\newcommand{\Tr}{\mathrm{Tr}}
\newcommand{\dbar}{d_{\bar{}}}

\title{Five-point genus-zero Arnold sewing}
\author{Raeez Lorgat}
\date{}

\begin{document}
\maketitle

\begin{abstract}
The genus-zero five-point sewing problem is the formal completion of
$\bar M_{0,5}$ along its rational normal-crossing boundary. For an
evaluation-generated chiral $\Ainf$-algebra with nondegenerate node
pairing and trace-class genus-zero sewing sums, the boundary value of
the five-point correlator is the contraction of a three-point block
with a four-point block over the node state space. Pulling Arnold's
flat connection back along the local KZ map gives the arity-five
ordered bar differential; the Arnold three-term relation together
with the infinitesimal braid relations gives
$\dbar_5^2=0$~\cite{Arnold69,Kohno87}. Under the finite-window
square-zero model, or under Positselski's second-kind coderived
comparison in the curved CDG ambient, the degree-five completed
bar-cobar counit is the corresponding quasi-isomorphism or
coacyclic-cone comparison~\cite{LV12,Positselski11}.
\end{abstract}

\section{Statement}
\label{sec:n5-statement}

Let $\cA$ be a chiral $\Ainf$-algebra on a smooth projective curve
$X$, and let $\cA^{\mathrm{ec}}\subset\cA$ be its
evaluation-generated core. Fix an Arnold-local affine or formal
genus-zero collision screen in $X$, or take $X=\mathbb{P}^1$ with
the projective linear relations imposed. Assume the following
genus-zero data.
\begin{enumerate}[label=\textup{(\alph*)}]
\item The three-, four-, and five-point chiral correlators
$\Phi^{\mathrm{ec}}_n$ are defined on
$\Conf_n^{\mathrm{ord}}(X)$ and are compatible with genus-zero
sewing.
\item The node state space in $\cA^{\mathrm{ec}}$ carries a
nondegenerate homogeneous pairing $\eta$, with inverse matrix
$(\eta^{ab})$ on each finite energy slice.
\item The genus-zero state sums below converge in the chosen sewing
completion.
\item The arity-five finite-window bar-cobar comparison lies in the
square-zero total model, or in a curved CDG ambient with the
second-kind coderived comparison. In particular,
$\cA^{\mathrm{ec}}$ is treated as a strong completion tower whose
finite quotients lie in the proved finite-stage bar-cobar regime.
\end{enumerate}

Let $\mathscr S_5^{(0)}$ be the formal completion of $\bar M_{0,5}$
along its rational boundary. Its irreducible boundary divisors are
indexed by stable bipartitions
\[
\{S,S^c\}, \qquad |S|=2,\quad |S^c|=3,
\]
and their codimension-two intersections are the rational stable
trees of type $2+2+1$~\cite{DeligneM69,Knudsen83}. Genus-one trace
factors occur on the clutching boundary of $\overline M_{1,n}$, not
on $\bar M_{0,5}$.

\begin{theorem}[Five-point genus-zero sewing]
\label{thm:n5-genus-zero-sewing}
On the formal screen $\mathscr S_5^{(0)}$, the leading boundary
restriction of $\Phi^{\mathrm{ec}}_5$ is characterised by the
following genus-zero sewing data.
\begin{enumerate}[label=(\roman*)]
\item \textbf{Rational separating factorisation.} For each boundary
divisor $D_{\{S,S^c\}}$ with $|S|=2$ and $|S^c|=3$,
\[
\Phi^{\mathrm{ec}}_5\big|_{D_{\{S,S^c\}}}
\;=\;
\sum_{a,b}
\Phi^{\mathrm{ec}}_{|S|+1}\bigl(\xi_S;\, e_a\bigr)\,
\eta^{ab}\,
\Phi^{\mathrm{ec}}_{|S^c|+1}\bigl(e_b;\, \xi_{S^c}\bigr),
\]
where $\{e_a\}$ is a homogeneous basis of the
evaluation-generated node sectors and $\eta^{ab}$ is the inverse
node-pairing matrix on the relevant finite energy slice.

\item \textbf{KZ--Arnold pullback.} On the chosen Arnold-local
genus-zero screen, the ordered bar differential of
$\barB^{\mathrm{ord}}_{\le 5}(\cA^{\mathrm{ec}})$ satisfies
\[
\dbar_{\!5}=\mathrm{KZ}^*\!\bigl(\nabla_{\mathrm{Arnold}}^{(5)}\bigr).
\]
Here $\nabla_{\mathrm{Arnold}}^{(5)}$ is the five-point Arnold
connection on the ordered configuration chart. Its curvature
vanishes by Arnold's three-term relation
\[
\eta_{ij}\wedge\eta_{jk}
+\eta_{jk}\wedge\eta_{ki}
+\eta_{ki}\wedge\eta_{ij}=0,
\qquad 1\le i<j<k\le 5,
\]
with $\eta_{ij}=d\log(z_i-z_j)$. The pullback identity gives
$\dbar_{\!5}^{\,2}=0$ after the infinitesimal braid relations for
the OPE residue operators are imposed. No global positive-genus or
arbitrary-curve Arnold-flatness statement is asserted here.

\item \textbf{Arity-five Maurer--Cartan component.} The sewing
twisting morphism
\[
\tau^{\mathrm{sew},5}
\;\in\;
\MC\bigl(\Conv(\barB^{\mathrm{ch}}_{\le 5}(\cA^{\mathrm{ec}}),\,
\cA^{\mathrm{ec}})\bigr)
\]
restricts to the lower arity sewing morphism through arity four
and has degree-five component given by the factorisation in
\textup{(i)}. The Maurer--Cartan equation in arity five is
equivalent to the KZ--Arnold flatness package in \textup{(ii)}:
Arnold's three-term relation, the infinitesimal braid relations for
the residue operators, and chiral OPE compatibility on
codimension-two rational strata.

\item \textbf{Completed counit on $\cA^{\mathrm{ec}}$.} The completed
arity-five counit
\[
\widehat\epsilon_5\colon
\widehat\Omega^{\mathrm{ch}}\bigl(\barB^{\mathrm{ch}}_{\le 5}
(\cA^{\mathrm{ec}})\bigr)
\;\xrightarrow{\;\sim\;}\;
\cA^{\mathrm{ec}}_{\le 5}
\]
is a quasi-isomorphism in the square-zero total model. In the
curved CDG ambient it is the corresponding coacyclic-cone
comparison in the second-kind coderived category. Its degree-five
component is the cobar dual of $\tau^{\mathrm{sew},5}$.
\end{enumerate}
\end{theorem}

\begin{remark}[Evaluation-generated core]
\label{rem:n5-evaluation-core}
Theorem~\ref{thm:n5-genus-zero-sewing} makes no full-category
assertion. Extension beyond $\cA^{\mathrm{ec}}$ requires trace-class
OPE-growth control for non-evaluation sectors and the corresponding
full-category finite-stage bar-cobar comparison.
\end{remark}

\begin{remark}[Elliptic trace and genus-one clutching]
\label{rem:n5-elliptic-trace-separate}
The elliptic trace expression
\[
\Tr_{\cA^{\mathrm{ec}}}^{(1)}
\bigl(\Phi^{\mathrm{ec}}_3(\xi_3,\xi_4,\xi_5;e_a)
q^{L_0-c/24}\bigr)
\]
belongs to the boundary of a genus-one clutching surface, for
example a divisor in \(\overline M_{1,n}\) or in a universal
genus-one family. The rational boundary of \(\bar M_{0,5}\)
contains only genus-zero separating strata.
\end{remark}

\section{Proof}
\label{sec:n5-proof}

\begin{proof}[Proof of Theorem~\ref{thm:n5-genus-zero-sewing}]
The divisor $D_{\{S,S^c\}}$ has a formal normal coordinate $t$.
Expanding $\Phi^{\mathrm{ec}}_5$ in $t$ gives the OPE expansion
through the node. The constant term is the contraction of the
three-point correlator on one rational component with the four-point
correlator on the other rational component. Nondegeneracy of
$\eta$ gives the inverse matrix $(\eta^{ab})$, and the sewing
completion hypothesis gives convergence of the displayed state sum.
This proves \textup{(i)}. Compatibility on intersections of two
boundary divisors is associativity of the genus-zero sewing
operation, expressed on the two rational stable trees of type
$2+2+1$.

The ordered bar differential records the residues of the chiral OPE
along the Arnold hyperplanes on the chosen genus-zero collision
screen. Under the local KZ--Arnold map this is the pullback of
$\nabla_{\mathrm{Arnold}}^{(5)}$, giving \textup{(ii)}. Arnold's
three-term relation and the infinitesimal braid relations give
\[
\bigl(\nabla^{(5)}_{\mathrm{Arnold}}\bigr)^2=0,
\]
and hence $\dbar_5^2=0$ on the pulled-back ordered bar complex.

The degree-five component of $\tau^{\mathrm{sew},5}$ is the
boundary factorisation in \textup{(i)}, extended over the interior
by the pulled-back KZ differential. The identity
\[
d\tau+\frac12[\tau,\tau]=0
\]
in
$\Conv(\barB^{\mathrm{ch}}_{\le 5}(\cA^{\mathrm{ec}}),\cA^{\mathrm{ec}})$,
is the arity-five KZ--Arnold flatness package, with the bracket
term recording the two-level rational boundary strata. This proves
\textup{(iii)}.

The finite-window bar-cobar hypothesis applies to the arity-five
truncation of the strong completion tower $\cA^{\mathrm{ec}}$.
Therefore the compatible finite-stage counits assemble to
$\widehat\epsilon_5$. In the square-zero total model this assembled
counit is a quasi-isomorphism; in the curved CDG ambient the cone is
coacyclic under the stated second-kind hypotheses. The construction
identifies the degree-five component with the cobar dual of
$\tau^{\mathrm{sew},5}$, giving \textup{(iv)}.
\end{proof}

\begin{remark}[Rational and elliptic boundary strata]
\label{rem:n5-rational-elliptic}
The boundary of $\bar M_{0,4}$ is a three-point set. The boundary
screen of $\bar M_{0,5}$ consists of ten rational divisors meeting
along codimension-two rational strata. Elliptic trace factors enter
through the clutching morphisms for \(\overline M_{1,n}\), where
the node carries a trace over the state space. Genus-zero boundary
screens use inverse-pairing contractions.
\end{remark}

\begin{thebibliography}{99}

\bibitem{Arnold69}
V.~I. Arnold,
\emph{The cohomology ring of the colored braid group},
Mat. Zametki \textbf{5} (1969), 227--231.

\bibitem{DeligneM69}
P.~Deligne and D.~Mumford,
\emph{The irreducibility of the space of curves of given genus},
Publ. Math. IH\'ES \textbf{36} (1969), 75--109.

\bibitem{Knudsen83}
F.~F. Knudsen,
\emph{The projectivity of the moduli space of stable curves, II:
The stacks $M_{g,n}$},
Math. Scand. \textbf{52} (1983), 161--199.

\bibitem{Kohno87}
T.~Kohno,
\emph{Monodromy representations of braid groups and Yang--Baxter
equations},
Ann. Inst. Fourier (Grenoble) \textbf{37} (1987), no.~4, 139--160.

\bibitem{LV12}
J.-L.~Loday and B.~Vallette,
\emph{Algebraic Operads},
Grundlehren der mathematischen Wissenschaften, vol.~346,
Springer, 2012.

\bibitem{Positselski11}
L.~Positselski,
\emph{Two kinds of derived categories, Koszul duality, and
comodule-contramodule correspondence},
Mem. Amer. Math. Soc. \textbf{212} (2011), no.~996.

\end{thebibliography}

\end{document}
